\documentclass{article}
\usepackage[margin=0.8in]{geometry}
\usepackage{kotex}


\begin{document}

\section{자기소개서}

\subsection{자기소개 (1600 byte)}

저는 로봇 인공지능 연구를 위해 2026년 가을학기 Virginia Tech 기계공학 박사과정에 입학하는 지원자 이현빈입니다. 2022년 및 2024년 경희대학교에서 기계공학 학/석사 학위를 취득하며 데이터 기반 기계학습으로 복잡한 물리 시스템을 모델링하는 방법을 탐구해 왔고, 그 결과 석사과정 중 1저자 논문 1편을 포함한 총 2편의 SCIE 저널 논문 게재 및 국제 학회 포함 총 5회의 1저자 컨퍼런스 발표 성과를 남겼습니다. 이후 저의 학문적 호기심은 모델링을 넘어 제어로까지 확장되었고, 석사 졸업 직전 KAIST 김재철 AI대학원 Robot Learning 연구실에서의 경험을 통해 학습 기반 로봇 제어에 큰 관심을 두게 되었습니다. 석사 졸업 이후에는 또 다른 삶의 방향을 탐색하기 위해 공학 학위자로서는 이례적으로 패션업계에 종사하며 비주얼 디렉팅, 모델, 커뮤니케이션 업무를 통해 새로운 경험을 쌓았고, 동시에 퇴근 후 경희대학교 소속으로 후속 연구에 몰두하며 약 1년간 학계와 예술계에 걸친 경험을 이어갔습니다. 이 기간 동안 제도적 보상 없이도 이어지는 공학 연구에 대한 지속된 열망을 확인했고, 이후 미국 박사과정에 지원하여 총 2개 학교의 전액 장학금 오퍼를 받게 되었습니다. 현재는 경희대학교에서의 후속 연구로 IEEE/ASME T-MECH 1저자 논문을 포함한 총 2편의 로봇 분야 논문을 SCIE 저널에 추가 제출한 상태이며, 연구를 로봇의 접촉 작업 제어를 위한 물리 기반 인공지능 모델로 확장해 나가고 있습니다. 지금까지 저의 여정은 경계를 드나드는 탐험의 연속이었습니다. 대학원에서는 연구 분야의 경계를, 졸업 후에는 과학과 예술의 경계를 넘으며 현재의 방향을 확립했습니다. 앞으로 저는 급변하는 로봇 및 인공지능 기술의 흐름 속에서, 확고한 방향성과 함께 인공지능과 로봇공학을 탐험하며 인간 수준의 로봇 제어를 실현하기 위한 연구를 지속해 나갈 것입니다.


\subsection{지원동기 (1600 byte)}

저는 학/석사 재학 시절 모교인 경희대학교로부터 총 4,500만 원가량의 재정적 지원을 받으며 학위과정을 마쳤습니다. 덕분에 당시 저는 연구에서 재정적 지원의 중요성을 크게 느끼지 못했습니다. 그러나 졸업 이후 별도의 지원 없이 생업과 후속 연구를 병행하면서, 저는 시간적 압박과 커리어 불일치로 인해 처음으로 연구 지속의 어려움을 실감하였습니다. 이를 해결하기 위해 저는 프리랜서 개발자로 생활비를 충당하기도 하고, 제 사정을 이해해 주는 스타트업과 함께 일하기도 했으며, 모교 지도교수님께 연구실의 컴퓨팅 장비 이용을 허락받아 연구 비용을 절감하기도 했습니다. 그 과정 속에서 저는 어려움 속에서도 지속되는 스스로의 연구 의지를 명확히 확인했고, 동시에 안정적인 경제적 지원이 연구에 얼마나 중요한지도 절실히 느꼈습니다. 또한 이 시기 저는 연구실 후배들의 멘토링과 연구실 홈페이지(GitHub) 관리도 함께 수행하며, 유학을 위한 성과 독식이 아닌 제 지식과 능력을 공동체에 기여하는 방향으로 연구에 임하였습니다. 결국 주변의 여러 도움과 노력에 힘입어 저는 미국 박사과정 입시와 후속 연구를 모두 성공적으로 마무리하고, 이제는 새로운 여정을 떠나게 되었습니다. 제가 선택한 Virginia Tech의 박사과정은 매년 여름을 제외한 9개월, 총 4년간의 재정지원을 보장합니다. 일반적으로 박사과정이 5년가량 소요된다는 점을 고려할 때, 일주학술문화재단의 지원은 향후 저의 박사과정 연구 안정성 확보에 큰 도움이 될 것으로 생각합니다. 비록 제가 구조적 지원이 없는 환경에서도 2년간 독립적으로 연구를 지속하며 성과를 내왔지만, 경제적 여건 때문에 온전히 연구에 집중하지 못한 시기 또한 길었다고 생각합니다. 일주재단의 지원은 제가 이러한 부담을 덜고, 과학계 공동체와 기술 발전에 기여하는 독립 연구자로 성장하는 데 크나큰 보탬이 될 것입니다.


\subsection{일주재단과 설립자에 대한 생각 (1600 byte)}

대한민국은 폐허 속에서 기술과 인적 자원을 바탕으로 선진국에 이르렀습니다. 이 과정에서 고 이임용 회장께서는 건국 초기 태광그룹을 이끌며 국가 산업 발전에 기여했을 뿐 아니라, 미래 인재 양성을 위한 기반 마련에도 많은 힘을 써주셨습니다. 특히 그룹의 장학 지원과 세화 중고교 운영이 현재 수십 년 넘게 이어지고 있다는 점에서, 이임용 회장께서 얼마나 기업의 사회 환원과 인재 교육의 중요성을 깊이 강조하셨던 분이셨는지 알 수 있습니다. 일주 선생께서 살아온 시대는 모두가 생존을 위해 분투하던 시기였습니다. 당장 오늘을 걱정해야 하는 환경에서 기업을 일구고, 정치적 혼란 속에서도 올곧음을 지키며 후대를 위한 기반까지 마련하고자 했던 그 마음을 제가 온전히 이해할 수는 없을 것입니다. 다만, 후대의 배움을 위해 힘쓰고자 했던 그 의지만큼은 분명하게 전해집니다. 저는 학업과 연구를 이어가며, 한 사람의 성장은 개인의 의지만이 아닌 구조적 지원 위에서 가능하다는 사실을 절실히 느꼈습니다. 그런 점에서 일주재단이 이어가고 있는 장학의 의미는 더욱 무겁게 다가옵니다. 현재 우리는 산업화 세대가 다진 기반 위에서 많은 기회를 누리고 있지만, 그 유산을 잇고 발전시켜야 할 책임도 함께 지니고 있습니다. 풍요 속에서도 발전의 성장통은 여러 사회 문제로 이어졌고, 고령화와 저출산으로 국가 경쟁력은 저하되고 있으며, 주요 경쟁국들은 자본과 인구를 앞세워 우리나라를 빠르게 추격하고 있습니다. 이런 현실에서 자원과 영토가 한정된 우리나라에게 인재와 기술은 변함없이 가장 중요한 자산입니다. 젊은 공학자로서, 저는 박사 유학을 통해 제 분야에서 깊이 있는 성과를 쌓는 데 그치지 않고, 졸업 후 제가 받은 배움과 경험을 우리 사회와 후속 세대에 나누며 주어진 책임을 다하고자 합니다. 저는 그것이 이임용 회장님께서 남기신 뜻에 응답하는 길이라 믿습니다.

\pagebreak

\section{연구목표 및 학위 취득 후 계획 (3500 byte)}

로봇이 실생활에서 도움을 주기 위해서는, 주어진 작업을 이해하고 환경 및 대상 물체와 올바르게 상호작용할 수 있어야 합니다. 이를 위해 로봇공학은 주변 환경의 정확한 인식, 대상에 안전하고 효율적으로 접근하는 경로 계획, 그리고 실제로 물체와 접촉하여 작업을 수행하는 조작 제어의 세 축을 중심으로 발전해 왔습니다. 저의 이전 연구는 이 중 인식과 깊게 맞닿아 있고, 박사과정을 통해 접촉 조작 제어 분야로 기존 연구를 확장하고자 합니다. 

최근 인공지능 및 기계학습을 로봇에 접목하는 연구가 활발히 진행되며, 기존엔 어려웠던 복잡한 환경 인식 및 제어가 점점 가능해지고 있습니다. 이처럼 로봇공학은 많은 발전을 이뤘지만, 실제 적용을 위해서는 아직 난관이 많습니다. 저의 목표인 정확하고 효율적인 접촉 조작을 위해서는, 다양한 로봇의 접촉 동역학 상태를 통합하여 제어 모델에 반영해야 합니다. 그러나 접촉에 중요한 외력/촉각과 같은 정보는 위치, 모터 토크 등 고유수용감각과 달리 획득을 위해 추가 센서 장착 또는 정교한 상태 추정이 필요하며, 이는 현재 조작 제어의 난관 중 하나로 남아 있습니다. 비용, 공간 등 요인으로 센서 장착이 어렵거나, 상태 추정이 어려운 시스템도 많기 때문입니다. 여기서 저의 기존 연구들은 물리 지식을 반영한 기계학습 기법으로 접촉과 진동이 빈번한 환경에서의 동역학 상태를 빠르고 정교하게 예측합니다. 가장 최근 연구(대표 저서)에서는, 확률모델과 전이 학습을 통해 추가 센서 없이 고유수용감각만으로 고주파 외력을 정교하게 예측하여, 고진동 환경에서의 데이터 기반 상태 예측을 가능하게 하였습니다. 이러한 경험을 토대로, 저는 향후 박사과정에서 로봇의 접촉 상태를 센서 없이 정확히 추정하고 제어에 통합하는 방향으로 더 정밀한 조작을 가능하게 할 계획입니다.

제어 또한 많은 기술적 요구사항을 만족해야 합니다. 기존의 고전적인 제어 방법은 해석 가능하고 안정적이지만 복잡한 작업 지시에 한계가 있고, 학습 기반 방법은 유연하고 적응성이 높지만 다양한 환경에서의 범용성, 보상 설계의 어려움, 내부 작동 원리 해석 등 해결되지 못한 문제가 많습니다. 최근에는 대규모 사전 학습 모델, 모방 및 강화 학습, 시뮬레이션-현실 간 전이, 고전 및 학습 제어의 하이브리드 방법을 통해 이러한 한계를 점차 극복해 나가고 있으며, ChatGPT와 같은 거대 언어 모델(LLM)을 제어에 통합하는 시도 또한 적극 연구되고 있습니다. Virginia Tech에서 저를 지도해주실 Simon Stepputtis 교수님은 이 중 LLM이 가진 범용 지식과 자연어 처리 능력을 로봇 제어에 통합하는 연구를 수행하고 계시며, 고수준 LLM 경로 계획과 저수준 강화 학습의 결합, LLM을 통한 작업의 계층적 분해 및 제어, 그리고 모듈 방식의 해석 및 전이 가능한 정책 모델 등을 제안하며 분야를 선도하고 계십니다. 저는 앞선 주제와 더불어, Simon 교수님의 방향에 저의 동역학 기반 지식과 물리 시뮬레이션 경험을 융합하여, LLM을 통한 제어가 물리적으로 올바르고 안전할 수 있도록 제어 모델 연구 또한 함께 수행할 예정입니다.

학위 취득 후에는 로봇공학과 인공지능의 교차점에서 장기적인 연구를 이어가는 연구자가 되고자 합니다. 현재로썬 학계와 연구 중심 산업 조직 모두 가능성을 열어두고 있지만, 제가 얻은 배움을 나누는 데 직접 기여할 수 있다는 점에서 학계로 더 무게를 두고 있습니다. 제가 이 길에 오르기까지, 감사하게도 많은 분의 지도와 지원이 있었습니다. 특히 석사 지도교수님께는 5년이 넘는 가르침을 받아왔습니다. 그럼에도 졸업 이후 제가 감사의 마음으로 커피 한 잔이라도 사서 찾아뵈면, 교수님께선 제가 열심히 연구하고 성장하는 것이 보은이니, 매번 다음엔 빈손으로 오라고 말씀하십니다. 저는 제가 이번 유학을 통해 더 나은 연구자로 성장하고, 제가 받은 가르침을 미래에 우리나라 연구자들과 사회에 다시 나누며 기여하는 것이, 이러한 은덕에 보답하는 길이라 생각합니다. 일주 이임용 회장님께서 남기신 ‘나무는 숲과 함께 자라야 한다’라는 생각 또한, 이러한 감사와 보은의 마음이 아닐까 하는 사견과 함께 글을 마칩니다. 감사합니다.

\end{document}